\chapter{System Requirements and Feasibility Analysis}
\label{ch:requirements}

\section{Introduction}
Engineering an operating system-level voice intelligence and automation framework demands a comprehensive requirements specification and a rigorous multi-dimensional feasibility study. Because PRIVACY68 executes arbitrary system automation commands, interacts with hardware audio transducers, and orchestrates diverse foundation models, its operational envelope must be strictly demarcated. 

This chapter specifies the system requirements for PRIVACY68, categorizing them into functional requirements (FR), non-functional requirements (NFR), hardware profiles, and software dependencies. It provides a technical justification of the chosen technology stack and delivers an in-depth feasibility analysis covering technical, operational, and economic dimensions, including a quantitative Total Cost of Ownership (TCO) comparison.

\section{System Requirements Specification}

\subsection{Functional Requirements (FR)}
The functional requirements define the explicit behavioral capabilities, inputs, processing operations, and outputs of PRIVACY68. Table \ref{tab:functional_requirements} delineates the fifteen primary functional requirements of the system.

\begin{table}[htbp]
\centering
\scriptsize
\caption{Comprehensive Functional Requirements Specification}
\label{tab:functional_requirements}
\begin{tabularx}{\textwidth}{|l|p{2.8cm}|p{3.2cm}|X|c|}
\hline
\textbf{Req ID} & \textbf{Requirement Name} & \textbf{Input / Trigger} & \textbf{System Processing \& Output} & \textbf{Priority} \\ \hline
\textbf{FR-01} & Hardware-Isolated PTT Recording & User clicks HUD mic or holds PTT hotkey & Opens PortAudio 16kHz mono stream; flushes PCM buffer to RAM; physically closes device on release. & High \\ \hline
\textbf{FR-02} & Dynamic Peak Normalization & Raw 16kHz floating point PCM buffer $\mathbf{x}$ & Computes peak amplitude $P = \max |\mathbf{x}|$; scales buffer by $(\mathbf{x}/P) \times 0.9$ if $P > 0.005$; outputs scaled PCM. & High \\ \hline
\textbf{FR-03} & Primary Online Speech Recognition & Normalized 16kHz audio buffer & Dispatches audio to Google Web Speech API; parses JSON transcript within 100ms; falls back to FR-04 upon network failure. & High \\ \hline
\textbf{FR-04} & Offline Neural Speech Fallback & Normalized 16kHz audio buffer & Initializes local Faster-Whisper (int8 CTranslate2) engine; computes beam search transcription locally. & High \\ \hline
\textbf{FR-05} & Lexical Preprocessing \& Autocorrect & Raw transcribed string & Strips enclosing quotes/backticks; removes conversational fillers; corrects domain hotwords (e.g., whatsapp, browser). & High \\ \hline
\textbf{FR-06} & Fast-Path Intent Dispatch & Preprocessed command string & Matches commands against built-in dictionary (media keys, volume, launch app, web URLs); dispatches in $<10$ms. & High \\ \hline
\textbf{FR-07} & Multi-Provider Model Cascading & Unhandled natural language prompt & Queries LiteLLM gateway with primary model (Gemini 2.5 Flash); automatically cascades to Groq LPU or Ollama on HTTP 429/503. & High \\ \hline
\textbf{FR-08} & Autonomous Script Synthesis & LLM response payload & Extracts fenced Python automation code; validates basic syntax; passes to execution sandbox. & High \\ \hline
\textbf{FR-09} & Subprocess Execution Sandbox & Validated Python script & Spawns isolated worker process (\texttt{subprocess.run}); enforces 30s timeout; captures stdout and stderr streams. & High \\ \hline
\textbf{FR-10} & Self-Healing Code Refinement & Execution error or non-zero exit code & Combines original prompt, failed code, and stderr traceback into new LLM prompt; re-synthesizes corrected script (max 3 tries). & High \\ \hline
\textbf{FR-11} & Workflow Skill Crystallization & Successfully executed script & Extracts action keywords; appends verified script to \texttt{skills.json}; provides instant future recall without LLM query. & Medium \\ \hline
\textbf{FR-12} & 3-Tier Security Interception & Synthesized script or shell action & Analyzes AST for high-risk operations (file deletion, git push, process kill); pauses execution and issues 15s confirmation prompt. & High \\ \hline
\textbf{FR-13} & Ergonomic HUD \& System Tray & User interaction / OS event & Renders glassmorphic floating HUD via Edge WebView2; provides system tray icon with real-time status and voice trigger. & Medium \\ \hline
\textbf{FR-14} & Mobile Wi-Fi Web Remote & Mobile browser HTTP request & Serves authenticated responsive HTML5 remote over local LAN; validates 6-character PIN; executes remote commands. & Medium \\ \hline
\textbf{FR-15} & Encrypted Telegram Teleoperation & Telegram bot message / voice note & Transcribes voice notes via Faster-Whisper; executes workstation commands; transmits desktop screenshot confirmations. & Low \\ \hline
\end{tabularx}
\end{table}

\subsection{Non-Functional Requirements (NFR)}
Non-functional requirements specify the operational quality attributes, performance envelopes, security standards, and ergonomic criteria governing PRIVACY68, as summarized in Table \ref{tab:non_functional_requirements}.

\begin{table}[htbp]
\centering
\small
\caption{Non-Functional Requirements Specification}
\label{tab:non_functional_requirements}
\begin{tabularx}{\textwidth}{|l|p{3.2cm}|X|}
\hline
\textbf{NFR ID} & \textbf{Quality Attribute} & \textbf{Specification Envelope} \\ \hline
\textbf{NFR-01} & \textbf{End-to-End Latency} & Fast-path commands must execute within $\le 100\text{ms}$. Full LLM agentic code synthesis via cloud LPU must complete within $\le 1,200\text{ms}$. \\ \hline
\textbf{NFR-02} & \textbf{Zero Idle Surveillance} & Ingestion microphone must remain closed at the driver level during idle states. Background acoustic buffer leak probability must be $0.0000$. \\ \hline
\textbf{NFR-03} & \textbf{Workstation Safety} & Zero destructive shell commands (directory wiping, process termination, branch deletion) shall execute without explicit human confirmation. \\ \hline
\textbf{NFR-04} & \textbf{Fault Tolerance} & System must survive cloud API rate limits, model deprecations, and network disconnections via automatic multi-tier fallback without crashing. \\ \hline
\textbf{NFR-05} & \textbf{Resource Footprint} & Idle RAM consumption shall not exceed $90\text{MB}$. Peak CPU utilization during speech capture and peak normalization shall remain below $5\%$. \\ \hline
\textbf{NFR-06} & \textbf{Speech Robustness} & Normalization must maintain Word Error Rates (WER) below $10\%$ across whispered speech and background office noise environments. \\ \hline
\textbf{NFR-07} & \textbf{Local Portability} & System must deploy as a standalone, zero-dependency executable installer on Windows 10 and 11 without requiring pre-installed Python. \\ \hline
\textbf{NFR-08} & \textbf{Auditability} & All security-intercepted operations, user confirmations, and execution errors must be permanently logged in a local audit journal. \\ \hline
\textbf{NFR-09} & \textbf{Visual Ergonomics} & Desktop HUD overlay must support transparency, draggable positioning, click-through capability, and non-blocking multi-window workflows. \\ \hline
\textbf{NFR-10} & \textbf{Remote Security} & Mobile remote access over LAN must enforce ephemeral 6-character PIN authentication with 401 Unauthorized lockout protection. \\ \hline
\end{tabularx}
\end{table}

\subsection{Hardware Requirements and Hardware Sizing Profiles}
PRIVACY68 is engineered to run efficiently across three distinct workstation tiers, ranging from standard corporate thin clients to high-performance local AI workstations, as detailed in Table \ref{tab:hardware_profiles}.

\begin{table}[htbp]
\centering
\small
\caption{Workstation Hardware Sizing Profiles}
\label{tab:hardware_profiles}
\begin{tabularx}{\textwidth}{|p{2.8cm}|X|X|X|}
\hline
\textbf{Component} & \textbf{Profile A: Thin Client / Office PC} & \textbf{Profile B: Developer Workstation} & \textbf{Profile C: Local AI Powerhouse} \\ \hline
\textbf{Target User} & Knowledge worker, enterprise productivity & Software engineer, power user & Air-gapped researcher, privacy purist \\ \hline
\textbf{Processor (CPU)} & Intel Core i3 / AMD Ryzen 3 (2.0 GHz, 2 Cores) & Intel Core i5/i7 11th Gen+ / AMD Ryzen 5000 (4--8 Cores) & Intel Core i9 / AMD Ryzen 9 / Threadripper (16+ Cores) \\ \hline
\textbf{System RAM} & 4 GB DDR4 & 16 GB DDR4/DDR5 & 32 GB--64 GB DDR5 \\ \hline
\textbf{Storage} & 1.5 GB free disk space (SATA HDD / SSD) & 5 GB free space (NVMe SSD) & 20 GB free space (NVMe PCIe 4.0 SSD) \\ \hline
\textbf{Audio Capture} & Integrated 3.5mm analog mic / webcam mic & USB desktop condenser microphone (16kHz / 48kHz) & Balanced XLR cardioid mic via USB audio interface \\ \hline
\textbf{Discrete GPU} & None (Intel UHD / Iris Xe) & Optional (NVIDIA RTX 3050/3060) & NVIDIA RTX 4080/4090 (16GB--24GB VRAM, CUDA) \\ \hline
\textbf{LLM Backend} & Cloud APIs (Gemini Flash, Groq LPU) & Cloud APIs + hybrid local fallback & 100\% Offline Local (Ollama: Llama-3.3-70B, Qwen2.5) \\ \hline
\textbf{ASR Backend} & Google Web Speech API & Google Web Speech + Whisper base & Local Faster-Whisper large-v3 via CUDA \\ \hline
\end{tabularx}
\end{table}

\subsubsection{Acoustic Sensor Specification}
The acoustic ingestion module interfaces with audio hardware via the PortAudio API. Supported audio input specifications include:
\begin{itemize}
    \item \textbf{Sampling Rate:} 16,000 Hz ($16\text{ kHz}$) mono capture. 16kHz provides a Nyquist frequency of 8kHz, completely covering the human vocal formant spectrum ($300\text{ Hz}$ to $3,400\text{ Hz}$) while minimizing memory consumption.
    \item \textbf{Bit Depth and Representation:} 32-bit IEEE 754 floating-point format (\texttt{float32}). Capturing directly in floating point provides $+1,500\text{ dB}$ of internal dynamic range, preventing digital arithmetic overflow prior to peak normalization.
    \item \textbf{Acoustic Buffer Block Size:} 1,280 samples ($80\text{ ms}$ chunks). This minimizes CPU interrupt frequency while ensuring sub-millisecond buffer aggregation.
\end{itemize}

\subsection{Software Requirements and Dependency Stack}
The system runtime environment is built on open-source, commercially permissive libraries:
\begin{itemize}
    \item \textbf{Host Operating System:} Microsoft Windows 10 (Build 19041+) or Windows 11 (64-bit).
    \item \textbf{Runtime Language:} Python 3.11.x or Python 3.12.x (64-bit).
    \item \textbf{WebView Rendering Subsystem:} Microsoft Edge WebView2 Runtime (evergreen distribution).
    \item \textbf{Key Python Dependencies:}
    \begin{itemize}
        \item \texttt{SpeechRecognition (>= 3.10.0)}: High-level ASR driver interface.
        \item \texttt{sounddevice (>= 0.4.6)}: Real-time C-PortAudio wrapper for hardware audio I/O.
        \item \texttt{numpy (>= 1.26.0, < 3.0)}: Vectorized numerical processing for dynamic peak scaling.
        \item \texttt{litellm (>= 1.50.0)}: Multi-provider Large Language Model abstraction gateway.
        \item \texttt{pywebview (>= 5.0.0)}: Lightweight Windows WebView2 container.
        \item \texttt{pystray (>= 0.19.5)}: Windows system tray notification area integration.
        \item \texttt{faster-whisper (>= 1.0.0)}: CTranslate2-quantized neural speech recognition engine.
        \item \texttt{requests (>= 2.31.0)}: Synchronous and asynchronous HTTP networking.
        \item \texttt{python-dotenv (>= 1.0.0)}: Dynamic environment variable management.
    \end{itemize}
    \item \textbf{Packaging Toolchain:} PyInstaller 6.x for binary synthesis; Inno Setup 6.x for standalone Windows installer generation.
\end{itemize}

\section{Technology Stack Justification}

\subsection{Microsoft Edge WebView2 vs. Electron.js}
Desktop GUI frameworks represent a critical architectural choice. Many contemporary AI assistants utilize Electron.js, bundling a dedicated Chromium browser instance and Node.js runtime with each application. Table \ref{tab:webview2_vs_electron} demonstrates why PRIVACY68 selects Microsoft Edge WebView2 via \texttt{pywebview}.

\begin{table}[htbp]
\centering
\small
\caption{Architectural Comparison: WebView2 vs. Electron.js}
\label{tab:webview2_vs_electron}
\begin{tabularx}{\textwidth}{|p{3.5cm}|X|X|}
\hline
\textbf{Evaluation Metric} & \textbf{Electron.js} & \textbf{Microsoft Edge WebView2 (PRIVACY68)} \\ \hline
\textbf{Idle RAM Consumption} & 220 MB -- 380 MB & \textbf{55 MB -- 85 MB (65\%--75\% Reduction)} \\ \hline
\textbf{Installer Package Size} & 110 MB -- 180 MB & \textbf{28 MB -- 35 MB (Standalone Binary)} \\ \hline
\textbf{Chromium Bundling} & Bundles redundant Chromium binary & Utilizes pre-installed Windows evergreen runtime \\ \hline
\textbf{Python Interoperability} & Requires complex IPC or WebSocket server & Native direct bi-directional JS-to-Python bridge \\ \hline
\textbf{Cold Launch Latency} & 1,400 ms -- 2,200 ms & \textbf{280 ms -- 420 ms} \\ \hline
\textbf{Security Patching} & Developer must recompile on CVE & Patched automatically by Windows Update \\ \hline
\end{tabularx}
\end{table}

By leveraging WebView2, PRIVACY68 retains full HTML5/CSS3 glassmorphic styling capabilities while minimizing system resource consumption.

\subsection{Sounddevice (PortAudio) vs. PyAudio}
For audio ingestion, \texttt{sounddevice} was selected over legacy \texttt{PyAudio}. \texttt{sounddevice} interfaces directly with the modern PortAudio v19 C library, providing zero-copy NumPy array buffer streaming. This eliminates memory serialization overhead, allowing instant vector peak normalization directly on incoming audio buffers.

\subsection{LiteLLM Unified Proxy Architecture}
Directly integrating individual foundation model SDKs (e.g., \texttt{google-generativeai}, \texttt{groq-python}, \texttt{openai}, \texttt{anthropic}) introduces severe code bloat, mismatched error-handling patterns, and fragile dependency trees. LiteLLM provides a standardized interface that translates vendor-specific REST protocols into a uniform completion format, enabling automatic intra-provider and cross-provider failover cascades in under 5 lines of configuration code.

\section{Feasibility Analysis}

\subsection{Technical Feasibility}
Technical feasibility assesses whether the project can be constructed, integrated, and maintained within existing algorithmic bounds and physical hardware limits.
\begin{enumerate}
    \item \textbf{Acoustic Processing Latency Feasibility:} Vector peak scaling on a 3-second audio buffer ($48,000$ float32 samples) requires computing $\max |\mathbf{x}|$ followed by scalar vector division. On a standard x86-64 CPU with AVX2 SIMD instructions, NumPy executes this computation in $0.118\text{ ms}$. This represents less than $0.15\%$ of a single audio frame, demonstrating that acoustic conditioning introduces zero perceptible latency.
    \item \textbf{Multi-Provider AI Resilience Feasibility:} LiteLLM abstracts the diverse wire-protocols of Google Gemini (REST JSON), Groq Cloud LPU (OpenAI-compatible), and local Ollama daemons. Dynamic intra-provider fallback guarantees that if a cloud endpoint returns HTTP 429 (Rate Limit) or HTTP 503 (Unavailable), alternative providers in the cascade are evaluated in memory without disrupting the user.
    \item \textbf{Sandboxed Subprocess Safety Feasibility:} Isolating generated Python scripts inside separate operating system processes (\texttt{subprocess.run}) with strict 30-second timeouts guarantees that infinite loops, blocking network calls, or memory leaks cannot crash or freeze the host desktop application.
    \item \textbf{Hardware Driver Control Feasibility:} By binding PortAudio stream instantiation strictly to the PTT button lifecycle, the physical audio device is energized only when required, achieving genuine $0.0000$ idle background exposure.
\end{enumerate}
Consequently, the technical feasibility of PRIVACY68 is definitively established.

\subsection{Operational Feasibility}
Operational feasibility assesses whether users, enterprise environments, and systems administrators will accept and adopt the software.
\begin{itemize}
    \item \textbf{Enterprise and Workplace Privacy Compliance:} Commercial voice assistants are banned in secure corporate and research environments due to always-on microphone surveillance. PRIVACY68 provides verifiable hardware isolation, complying with corporate confidentiality policies, GDPR, and CCPA standards.
    \item \textbf{Non-Intrusive Workflow Ergonomics:} The transparent, floating HUD overlay employs the Win32 \texttt{WS\_EX\_TRANSPARENT} click-through window attribute when minimized, ensuring that developers, gamers, and writers are never interrupted during intensive desktop tasks.
    \item \textbf{Zero-Configuration Cross-Device Teleoperation:} Serving an authenticated mobile remote over local Wi-Fi allows users to command their desktop from any smartphone browser without installing third-party mobile apps or creating proprietary cloud accounts.
\end{itemize}
Operational feasibility is exceptionally high.

\subsection{Economic Feasibility and Total Cost of Ownership (TCO)}
Economic feasibility evaluates the recurring financial obligations of the system against alternative commercial implementations.

\subsubsection{Total Cost of Ownership (TCO) Model}
Consider an active workstation user executing an average of $N = 300$ voice commands per day ($9,000$ commands per month). Table \ref{tab:tco_comparison} provides an economic cost breakdown comparing a commercial cloud-tethered voice agent stack (AWS Transcribe + OpenAI GPT-4o) against PRIVACY68.

\begin{table}[htbp]
\centering
\small
\caption{Monthly Total Cost of Ownership (TCO) Comparison ($N=9,000$ Requests/Month)}
\label{tab:tco_comparison}
\begin{tabularx}{\textwidth}{|p{3.8cm}|X|X|}
\hline
\textbf{Cost Component} & \textbf{Commercial Cloud Stack (AWS + GPT-4o)} & \textbf{PRIVACY68 Architecture} \\ \hline
\textbf{Speech Recognition (ASR)} & AWS Transcribe: $\$0.024 / \text{min} \times 450\text{ min} = \mathbf{\$10.80}$ & Google Web Speech API: $\mathbf{\$0.00}$ (Zero-Key Free Tier) \\ \hline
\textbf{LLM Reasoning Core} & OpenAI GPT-4o: $\$5.00 / 1\text{M tokens} \approx \mathbf{\$22.50}$ & Google Gemini 2.5 Flash / Groq Cloud LPU: $\mathbf{\$0.00}$ (Free Tier) \\ \hline
\textbf{Local Offline Fallback} & None (Fails when offline: requires cloud) & Local Faster-Whisper + Ollama: $\mathbf{\$0.00}$ (Free On-Device) \\ \hline
\textbf{Cloud Hosting Infrastructure} & Managed EC2 / Lambda: $\mathbf{\$15.00}$ & Local Desktop Execution: $\mathbf{\$0.00}$ (Local Subprocess) \\ \hline
\textbf{Total Monthly Cost} & \textbf{\$48.30 / month} & \textbf{\$0.00 / month} \\ \hline
\textbf{Annualized Total Cost} & \textbf{\$579.60 / year} & \textbf{\$0.00 / year (100\% Free)} \\ \hline
\end{tabularx}
\end{table}

PRIVACY68 incurs \textbf{\$0.00 in recurring infrastructure costs}, making it accessible to students, independent researchers, and cost-conscious enterprise organizations. Economic feasibility is absolute.
